\section{Related Work}

\subsection{Intra-city spatial inequality}

Research on spatial inequality has traditionally focused on differences across countries and regions. In the UK, this is reflected in longstanding debates on the North--South divide and broader patterns of uneven development, with persistent disparities in productivity, employment, income, and health (Martin et al., 2016; Bambra et al., 2014; McCann, 2016). Comparable macro-spatial inequalities have been documented elsewhere, including East--West divides in Germany and coastal--inland divergence in China (Iammarino et al., 2019; Wei \& Ye, 2009). This body of work has been influential in demonstrating that inequality is spatially structured, but its coarse geographical scale necessarily obscures variation within cities. More recent reviews emphasise that cities are not merely passive containers of inequality but active systems in which spatial form, infrastructure, and access jointly shape unequal outcomes within urban space (Sarkar, 2024).

Cities are not internally homogeneous socio-spatial units. Neighbourhoods within the same urban area can differ markedly in deprivation, housing quality, access to services, and health outcomes. The neighbourhood-effects literature shows that local context shapes individual trajectories through mechanisms such as peer interaction, institutional access, and information flows (Ellen \& Turner, 1997; Galster, 2012; van Ham et al., 2012). In parallel, the segregation literature demonstrates that the spatial separation of social groups is not merely descriptive but constitutive of inequality, as uneven access to opportunities becomes embedded in residential patterns (Massey \& Denton, 1988; Tammaru et al., 2016). From this perspective, intra-city inequality is not a residual of regional disparities but a distinct process operating at the urban scale.

Recent work has reinforced this shift in focus. Socio-economic segregation has been identified as a defining feature of contemporary European cities (Tammaru et al., 2016), while neighbourhood inequality has been shown to persist through durable social and institutional mechanisms (Sampson, 2012). Rising income inequality is also associated with increasing spatial segregation (Reardon \& Bischoff, 2011), suggesting that distributional and spatial processes are closely intertwined. More broadly, urban theory has long treated cities as key arenas through which inequality is produced and reproduced (Massey, 1974; Fainstein, 2014; Soja, 2010; Florida, 2017).

The local scale is not only empirically relevant but also theoretically salient. Relative deprivation is evaluated against proximate others (Runciman, 1966; Smith et al., 2012), and empirical evidence shows that well-being responds to local income comparisons (Luttmer, 2005). Recent UK-based work similarly highlights the perceptual salience of neighbourhood-scale inequality (Suss, 2023), while deprivation may also be experienced through daily activity spaces rather than residential location alone (Comber et al., 2022).

In the UK context, empirical studies of intra-urban inequality have often focused on specific dimensions such as housing markets, health disparities, or education outcomes, frequently at the level of individual cities or metropolitan regions(Bambra et al., 2014; Comber et al., 2022; Suss, 2023). While these studies provide valuable insights into localised processes of disadvantage, they are typically limited in geographical coverage and do not support systematic comparison across cities within a unified analytical framework.

\subsection{Spatial measures of inequality}

The Gini coefficient remains one of the most widely used summary measures of inequality due to its simplicity and comparability (Cowell, 2011). In the context of urban analysis, it provides a convenient scalar indicator of how unevenly a variable is distributed across neighbourhoods within a city. However, the conventional Gini coefficient is fundamentally aspatial, meaning that identical distributions can yield identical Gini values regardless of their spatial configuration.

This limitation is analytically significant. A city in which deprivation is concentrated in contiguous clusters differs fundamentally from one in which high- and low-deprivation neighbourhoods are locally interspersed, yet both may exhibit the same level of overall inequality. As a result, aspatial measures capture dispersion but not spatial structure. This distinction has long been recognised in the spatial inequality literature. Arbia (2001) separates the distribution of values from their spatial arrangement, while Rey and Janikas (2005) argue that inequality and spatial dependence should be analysed jointly. Bickenbach and Bode (2008) similarly demonstrate that conventional concentration measures are not well suited to capturing spatial phenomena.

Methodological responses have therefore focused on extending the Gini coefficient to incorporate spatial relationships. Rey and Smith (2013) propose a spatial decomposition that distinguishes inequality arising between neighbouring and non-neighbouring observations. Building on this, Panzera and Postiglione (2020) develop a Gini-correlation framework that quantifies the share of total inequality attributable to spatial clustering. Related approaches have also been applied in segregation analysis (T{\"u}rk \& {\"O}sth, 2023). These contributions establish that overall dispersion and spatial configuration are analytically separable dimensions and should not be collapsed into a single summary statistic. Compared with alternative spatial measures such as Moran's I, which captures spatial autocorrelation but not inequality magnitude, or segregation indices that are typically defined for discrete groups, the spatial Gini framework is particularly suited to continuous deprivation measures. It allows the joint assessment of dispersion and spatial arrangement within a single, comparable metric, making it well aligned with the objectives of intra-city inequality analysis.

Population weighting is a critical component of this framework. Small-area spatial units differ substantially in population size, and unweighted measures assign equal importance to areas rather than individuals. The theoretical basis for weighted inequality measurement is well established (Mookherjee \& Shorrocks, 1982), and its relevance to spatial analysis has been emphasised in recent work (Panzera \& Postiglione, 2020). For LSOA-based analysis, this implies that inequality should be measured on a person-weighted basis to reflect the distribution of deprivation experienced by residents rather than the distribution of values across arbitrary spatial units.

Despite these advances, an important gap remains. Much of the existing literature focuses on income, GDP, or other single-dimensional indicators and is often applied at regional or inter-municipal scales. Applications at the intra-city level are comparatively limited, and even fewer studies combine population weighting with spatial decomposition in a way that enables systematic comparison across cities. In particular, the integration of spatial inequality measures with multidimensional deprivation indicators at the fine neighbourhood scale remains underdeveloped.

\subsection{IMD and deprivation}

The Index of Multiple Deprivation (IMD) provides a suitable empirical basis for analysing intra-city inequality. The English Indices of Deprivation 2025 measure relative deprivation across 33,755 Lower-layer Super Output Areas (LSOAs), combining seven domains---income, employment, education, health, crime, barriers to housing and services, and living environment---into a composite indicator (UK Government, 2025). Its multidimensional structure and fine spatial resolution make it particularly well suited to capturing variation in deprivation within cities.

The conceptual foundation of deprivation lies in Townsend's (1987) definition of relative disadvantage, which emphasises observable exclusion from the living conditions and activities considered normal in society. This framework underpins the development of multidimensional measures that extend beyond income to capture broader forms of social and material disadvantage. In England, IMD has become the standard approach for identifying neighbourhood-level deprivation (McLennan et al., 2019) and is widely used in both research and policy contexts. Its consistency, official status, and policy relevance make it a robust basis for comparative urban analysis.

Existing research using IMD, however, remains fragmented in its analytical focus. A large share of the literature employs IMD descriptively to identify deprived areas or explain variation in outcomes such as health and educational attainment. Other work extends deprivation beyond the residential neighbourhood, for example, by incorporating activity spaces to capture exposure across the city (Comber et al., 2022). Related work further shows that inequality within cities is also reflected in mobility patterns and access to activity spaces, highlighting how exposure to opportunities varies across socio-economic groups beyond residential location (Gao et al., 2024). Alternative approaches use locally salient indicators, such as housing values, to measure neighbourhood inequality more directly (Suss, 2023), but these often lack comparability across cities.

As a result, the literature has developed along parallel rather than integrated lines. IMD-based studies tend to emphasise deprivation levels rather than internal inequality. Methodological work on spatial inequality tends to focus on economic indicators rather than multidimensional deprivation. UK-based studies of neighbourhood inequality are often limited to specific cities or indicators, reducing their comparability.

This fragmentation leaves a clear empirical gap. There is limited evidence on how multidimensional deprivation is distributed within cities, how unequal these distributions are across neighbourhoods, and to what extent inequality is spatially structured rather than locally mixed. This also implies that city-level deprivation cannot be treated as a single scalar condition: average deprivation, exposure to extreme deprivation, internal inequality, and spatial configuration may diverge substantially across cities. This highlights the need for a framework that integrates multidimensional deprivation with population-weighted and spatially decomposed measures of inequality across cities, which this study seeks to provide.
